\documentclass[letterpaper,10pt]{article}
\usepackage{fullpage}
\usepackage{amsmath,amssymb,mathtools}
\usepackage[T1]{fontenc}
\usepackage[utf8]{inputenc}
\usepackage{longtable,booktabs,array}
\usepackage[table]{xcolor}
\usepackage{hyperref}
\usepackage{natbib}
\usepackage{draftwatermark}
\bibliographystyle{plainnat}
\SetWatermarkScale{5}
\hypersetup{colorlinks=true, linkcolor={blue}, filecolor={Maroon}, citecolor={blue}, urlcolor={blue}}

\newcommand{\PSS}[2]{\mathrm{PSS}_{#1}^{\mathrm{#2}}\left(T;\boldsymbol{m}, \mathrm{age} \right)}
\newcommand{\Sepsis}[1]{\mathrm{Sepsis}^{\mathrm{#1}}\left(T;\boldsymbol{m}, \mathrm{age} \right)}
\newcommand{\SepticShock}[1]{\mathrm{SepsisShock}^{\mathrm{#1}}\left(T;\boldsymbol{m}, \mathrm{age} \right)}
\newcommand{\SI}{\mathrm{SI}\left(t\right)}
\newcommand{\m}[1]{m_{\mathrm{#1}}}
\newcommand{\ORS}{\mathrm{ORS}\left(t; \m{resp}\right)}
\newcommand{\IMV}{\mathrm{IMV}\left(t; \m{resp}\right)}
\newcommand{\PFR}{\mathrm{PFR}\left(t; \m{resp}\right)}
\newcommand{\SFR}{\mathrm{SFR}\left(t; \m{resp}\right)}

\NewDocumentCommand{\DS}{m O{\boldsymbol{m}, \mathrm{age}}}{\mathcal{D}[\mathrm{ #1 }]\left(t; #2 \right)}



\title{Operational Definition of the Phoenix Sepsis Score}
%\author{Peter E. DeWitt, Ph.D}
\date{\today}

\begin{document}
\maketitle
\renewcommand*\contentsname{Table of contents}
{
\hypersetup{linkcolor=}
%\setcounter{tocdepth}{3}
\tableofcontents
}

{
\hypersetup{linkcolor=}
\listoftables
}
\pagebreak

\section{Objective}\label{sec:objective}

Provide formal notation for the operational definition of the Phoenix
Sepsis Score as published in \citet{sanchezpinto_2024_development} and
\citet{schlapbach_2024_international} along with several exploratory aggregation schemes.  The different
scoring methods are based on when and how data is aggregated.  The differences
in the aggregation methods can dramatically change the Phoenix score and the
published cut point of two or more points for sepsis may not be applicable to
all aggregations.

\section{Notation}\label{sec:notation}

\subsection{Symbols}\label{sec:symbols}

This section is a quick guide to the symbols used throughout the document.
Detailed definitions for the functions will follow in
\S~\ref{sec:function-definitions}

\begin{longtable}{c p{2.1in} p{2.7in}}
\caption{Notation used in the operational definition of the Phoenix Sepsis Score.}
  \label{tab:symbols} \\

  \toprule
  \textbf{Symbol} & \textbf{Meaning} & \textbf{Example or Notes} \\
  \midrule
  \endfirsthead

  \multicolumn{3}{c}
  {{\bfseries \tablename\ \thetable{} -- continued from previous page}} \\
  \toprule
  \textbf{Symbol} & \textbf{Meaning} & \textbf{Example or Notes} \\
  \midrule
  \endhead

  \midrule
  \multicolumn{3}{r}{{Continued on next page}} \\
  \midrule
  \endfoot

  \bottomrule
  \endlastfoot

  \addlinespace[0.5em]
  \multicolumn{3}{l}{\textbf{Sets and Membership}} \\

    \rowcolor{gray!20}
    $\{\cdot\}$ &
    Set notation used to define collections of values satisfying a condition. &
    $\{x : x \le t\}$ denotes the set of values $x$ less than or equal to $t.$
    Also, $\boldsymbol{A} = \left\{1, 3, 5\right\}$ is read as "the set A
    consisting of the elements 1, 3, and 5."\\

    $\emptyset$ &
    Empty set. &
    Represents a set containing no elements. \\

    \rowcolor{gray!20}
    $\in$ &
    Set membership operator. &
    $t \in T$ indicates that time $t$ lies within the time window $T$. \\

    $\boldsymbol{A} \setminus \boldsymbol{B}$ &
    Remove elements of a set &
    Example $\left\{1, 2, 3\right\} \setminus \left\{2\right\} = \left\{1, 3\right\}$ \\

  \addlinespace[0.5em]
  \multicolumn{3}{l}{\textbf{Intervals}} \\

    \rowcolor{gray!20}
    $(a,b)$ &
    Open interval. &
    $X \in (a,b) \equiv a < X < b$. \\

    $[a,b)$ &
    Left-closed, right-open interval. &
    $X \in [a,b) \equiv a \le X < b$. \\

    \rowcolor{gray!20}
    $(a,b]$ &
    Left-open, right-closed interval. &
    $X \in (a,b] \equiv a < X \le b$. \\

    $[a,b]$ &
    Closed interval. &
    $X \in [a,b] \equiv a \le X \le b$. \\

  \addlinespace[0.5em]
  \multicolumn{3}{l}{\textbf{Definitions, Equality, and Logic}} \\

    \rowcolor{gray!20}
    $\coloneqq$ &
    Definition operator indicating that the quantity on the left is defined by
    the expression on the right. &
    $f(x) \coloneqq x^2$ defines the function $f(x)$ as $x^2$. \\

    $=$ &
    Equality relation. &
    $2 + 3 = 5$. \\

    \rowcolor{gray!20}
    $\equiv$ &
    Logical equivalence or identity. &
    $A \equiv B$ means that $A$ and $B$ are equivalent statements. \\

    $\vee$ &
    Logical OR operator. &
    $A \vee B$ is true if either $A$ or $B$ is true. \\

    \rowcolor{gray!20}
    $\wedge$ &
    Logical AND operator. &
    $A \wedge B$ is true if both $A$ and $B$ are true. \\

    $\forall$ &
    Universal quantifier (“for all”). &
    $\forall t \in T$ means for every $t$ in the time window $T$. \\

    \rowcolor{gray!20}
    $\exists$ &
    Existential quantifier (“there exists”). &
    $\exists\, t \in T$ means for some time $t$ in $T$. \\

  \addlinespace[0.5em]
  \multicolumn{3}{l}{\textbf{Arithmetic and Functional Operators}} \\

    \rowcolor{gray!20}
    $\sum$ &
    Summation operator. &
    $\sum_{i=1}^{4} A_i = A_1 + A_2 + A_3 + A_4$. \\

    $\max$ &
    Maximum operator. &
    $\max_{t \in T} f(t)$ evaluates $f(t)$ over all
    $t \in T$ and returns the maximum value. $\max \boldsymbol{A}$
    returns the largest value of the set $\boldsymbol{A}$, e.g., $\max\left\{1,
    2, 3\right\} = 3.$\\

    \rowcolor{gray!20}
    $\min$ &
    Minimum operator. &
    $\min_{t \in T} f(t)$ evaluates $f(t)$ over all
    $t \in T$ and returns the minimum value. $\min \boldsymbol{A}$
    returns the smallest value of the set $\boldsymbol{A}$, e.g., $\min\left\{1,
    2, 3\right\} = 1.$ \\

    $\times$ &
    Multiplication operator. &
    \\

    \rowcolor{gray!20}
    $\mathbb{I}\{A\}$ & Indicator function returning 1 if condition $A$ is true
    and 0 otherwise. & Defined in \eqref{eq:indicator}.\\

  \addlinespace[0.5em]
  \multicolumn{3}{l}{\textbf{Time and Indexing}} \\

    $t$ &
    Time measured in minutes relative to encounter start. &
    Negative values denote pre-encounter measurements,
    $t=0$ marks encounter start, and positive values occur during the encounter. \\

    \rowcolor{gray!20}
    $T$ &
    Time window measured in minutes. &
    $T=[0,1440)$ represents the first 24 hours following encounter start. \\

    $m_{resp}$  & Look-back window for respiratory related values & $[0, \infty)$\\
    \rowcolor{gray!20}
    $m_{vaso}$  & Look-back window for vasoactive medications & $[0, \infty)$\\
    $m_{bp}$   & Look-back window for blood pressure measurements & $[0, \infty)$\\
    \rowcolor{gray!20}
    $m_{lac}$   & Look-back window for lactate & $[0, \infty)$\\
    $m_{gcs}$   & Look-back window for GCS & $[0, \infty)$\\
    \rowcolor{gray!20}
    $m_{pupil}$ & Look-back window for Pupils & $[0, \infty)$\\
    $m_{coag}$  & Look-back window for coagulation related values & $[0, \infty)$\\
    $m_{endo}$  & Look-back window for endocrine related values & $[0, \infty)$\\
    $m_{immu}$  & Look-back window for immunologic related values & $[0, \infty)$\\
    $m_{hepatic}$  & Look-back window for hepatic related values & $[0, \infty)$\\
    $m_{renal}$  & Look-back window for renal related values & $[0, \infty)$\\

    \rowcolor{gray!20}
    $\boldsymbol{m}$ &
    Vector of look-back windows used for carry-forward operations. &
    \\

  \addlinespace[0.5em]
  \multicolumn{3}{l}{\textbf{Observation History and Carry-Forward}} \\

    \rowcolor{gray!20}
    $\tau_X(t)$ &
    Most recent observation time for variable $X$ at or before time $t$. &
    $\tau_X(t)$ identifies when $X$ was last measured prior to $t$.
    Defined in \eqref{eq:tau}.\\

    $\Delta_X(t)$ &
    Staleness of measurement for variable $X$. &
    $\Delta_X(t)$ is the elapsed time since $\tau_X(t)$.
    Defined in \eqref{eq:staleness}.\\

    \rowcolor{gray!20}
    $\mathcal{C}[X](t;m)$ &
    Carry-forward operator returning the most recent value of $X$
    within $m$ minutes of time $t$. &
    Returns the last observed value of $X$ if $\Delta_X(t) \le m$,
    otherwise $\mathrm{NA}$.
    Defined in \eqref{eq:carry-forward}.\\
    \\

    $\DS{x}[\cdot]$ & Scoring operator: the point contriubiton of scoring entity $x$ at time $t.$ &
    When the common argument list $\left(\boldsymbol{m}, \mathrm{age}\right)$
    is used, arguments not relevant to $x$ are ignored.
    \\

    \rowcolor{gray!20}
    $\mathrm{NA}$ &
    Missing, unknown, or undefined value. &
    Returned when a required measurement is unavailable. \\

\end{longtable}

Two points are especially important for reading the rest of the
document. First, $\coloneqq$ is \textbf{not} the same as $=$. $\coloneqq$ means
we are \emph{defining} something. Second, $\mathrm{NA}$ means
``missing/unknown/undefined'' (e.g., a lab value was not measured),
which is different from a numeric value like 0. By contrast,
$\emptyset$ is the \textbf{empty set} (no elements), which we use in
set-based definitions like $\max \emptyset = \mathrm{NA}$.

For arithmetic, when $\mathrm{NA}$ is involved, the outcome is
always $\mathrm{NA},$ e.g., $\mathrm{NA} + X \coloneqq \mathrm{NA}$ $\forall X.$
Table~\ref{tab:logic-with-na} explicitly defines how logical statements are
interpreted when
working with $\mathrm{NA}$.
When a value is missing, we avoid forcing it into ``true'' or ``false.''
The rules below show how missingness propagates through logical
statements. If either side of an ``and/or'' statement is unknown, the
whole statement may be unknown unless the other side already determines
the answer.

\begin{longtable}{ll}
  \caption{Logic with $\mathrm{NA}$.}\label{tab:logic-with-na} \\
\toprule
Statement & Value \\
\midrule
\endhead
\bottomrule
\endlastfoot
  \multicolumn{2}{l}{\textbf{`OR' statements}} \\
    \rowcolor{gray!20}
    $\mathrm{NA} \vee \mathrm{NA}$ & $\mathrm{NA}$ \\
    $\mathrm{NA} \vee \text{true}$ & true \\
    \rowcolor{gray!20}
    $\mathrm{NA} \vee \text{false}$ & $\mathrm{NA}$ \\
  \multicolumn{2}{l}{\textbf{`And' statements}} \\
    \rowcolor{gray!20}
    $\mathrm{NA} \wedge \mathrm{NA}$ & $\mathrm{NA}$ \\
    $\mathrm{NA} \wedge \text{true}$ & $\mathrm{NA}$ \\
    \rowcolor{gray!20}
    $\mathrm{NA} \wedge \text{false}$ & false \\
\end{longtable}

\subsection{Function Definitions}\label{sec:function-definitions}

This section defines how we turn time-stamped clinical data into a
single value at any time $t$. In real EHR data, measurements are taken
at irregular times. The definitions below formalize ``what is the most
recent value?'' and ``is it too old to trust?'' so that scoring is
reproducible.

For any time-indexed physiologic variable $X,$ we write $X(t)$ for the value
recorded at time $t$ where $t$ is the time, in minutes from the encounter start,
i.e., $t=0$ denotes encounter starts, $t=32$ is 32 minutes into the encounter,
and $t=-12$ is 12 minutes before encounter start. The negative times, represent
measurements recorded before encounter start, such as EMS data.

For any variable $X,$ define the last observation time prior to or at
$t$ as:

\begin{equation}
  \label{eq:tau}
  \tau_{X}(t) \coloneqq \max\left\{s \leq t : X(s) \text{ observed or reported}\right\}, \quad \max \emptyset = \mathrm{NA}.
\end{equation}

Define the staleness of the measurement:

\begin{equation}
  \label{eq:staleness}
  \Delta_{X}(t) \coloneqq
    \begin{cases}
      t - \tau_{X}(t) & \tau_X(t) \neq \mathrm{NA} \\
      \mathrm{NA} & \text{otherwise.}
    \end{cases}
\end{equation}

We define a staleness-limited carry-forward operator. Carry forward when
working with retrospective data, a look back limit when working with
prospective or real time data. For $m \geq 0$ minutes,

\begin{equation}
  \label{eq:carry-forward}
  \mathcal{C}\left[X\right](t;m)\coloneqq
    \begin{cases}
      X\left(\tau_{X}(t)\right) & \Delta_{X}(t) \leq m \\
      \mathrm{NA} & \text{otherwise.}
    \end{cases}
\end{equation}

We will use an indicator function that returns 1 or 0:
\begin{equation}
  \label{eq:indicator}
  \mathbb{I}\left\{A\right\} \coloneqq
    \begin{cases}
      1 & A \text{ is true} \\
      0 & A \text{ is false or } \mathrm{NA}.
    \end{cases}
\end{equation}

In words: $\tau_X(t)$ tells you when the last measurement happened,
$\Delta_X(t)$ tells you how old it is, and $\mathcal{C}[X](t;m)$
returns that last value if it is recent enough (within $m$ minutes);
otherwise it returns ``missing.'' This is how we prevent using values
that are too old to reflect current physiology. The indicator
$\mathbb{I}\{A\}$ then lets us convert logical statements (e.g.,
``lactate is at least 5'') into 0/1 points for scoring.

For potentially incomplete time-indexed function $f,$ extrema are evaluated over
the observed values of $f.$ Define:
\begin{equation}
  \mathcal{A}_{f}\left(T\right) \coloneqq \left\{ f \left(t\right) : t \in T, f\left(t\right) \text{ is available}\right\}.
\end{equation}
Then
\begin{equation}
  \max_{t \in T} f\left(t\right) \coloneqq
    \begin{cases}
      \max \mathcal{A}_{f}\left(T\right) & \mathcal{A}_{f}\left(T\right) \ne \emptyset \\
      \mathrm{NA} & \mathcal{A}_{f}\left(T\right) = \emptyset,
    \end{cases}
\end{equation}
and
\begin{equation}
  \min_{t \in T} f\left(t\right) \coloneqq
    \begin{cases}
      \min \mathcal{A}_{f}\left(T\right) & \mathcal{A}_{f}\left(T\right) \ne \emptyset \\
      \mathrm{NA} & \mathcal{A}_{f}\left(T\right) = \emptyset.
    \end{cases}
\end{equation}

\section{Defining Organ System Dysfunction Scores}\label{sec:defining-organ-system-dysfunction-scores}

We assess four organ systems: respiratory, cardiovascular,
neurologic, and coagulation. Each subsection below lists the EHR inputs
needed to compute that organ score, then shows how those inputs are
transformed into constructed variables and point assignments. The goal
is to make every step of the scoring reproducible: given the same
time-stamped data, you should arrive at the same score.

The listed patient, laboratory, observation, event, and medication fields are as
we expect them to appear in the electronic health record. These are inputs, not
judgments about clinical severity.  We define a domain for the variable. These
domains are intentionally broad (wider than clinically plausible values) to keep
data validation separate from scoring.

\subsection{Respiratory}\label{sec:resp}
\subsubsection{Data from the EHR}\label{sec:resp-ehr}

There are several values we expect to get from the Electronic Health Record
capturing oxygenation, $FiO_2$ (fraction of inspired oxygen), $SpO_2$ (pulse
oximetry oxygen saturation), and $PaO_2$ (arterial oxygen pressure).
Additionally, the presence of mechanical or noninvasive respiratory support.
Together, they allow us to quantify how well the patient is oxygenating and
whether that oxygenation depends on advanced support.

\begin{longtable}{lllll}
  \caption{EHR Inputs for Respiratory Organ Dysfunction Scoring}
  \label{tab:resp-ehr} \\

  \toprule
  \textbf{Variable} &
  \textbf{Description} &
  \textbf{Type} &
  \textbf{Units} &
  \textbf{Domain} \\
  \midrule
  \endhead

  \bottomrule
  \endlastfoot

  \rowcolor{gray!20}
  $FiO_2$ &
  Fraction of Inspired Oxygen &
  Real &
  unitless &
  $[0.21, 1.00]$ \\

  $SpO_2$ &
  Pulse oximetry oxygen saturation &
  Integer &
  unitless &
  $\{0, 1, \ldots, 100\}$ \\

  \rowcolor{gray!20}
  $PaO_2$ &
  Arterial oxygen pressure &
  Real &
  mmHg &
  $[0, \infty)$ \\

  $P_{aw}^{VENT}$ &
  Mean Airway Pressure, reported from a ventilator &
  Real &
  cm H$_2$O &
  $[0, \infty)$ \\

  \rowcolor{gray!20}
  $P_{aw}^{HFOV}$ &
  Mean Airway Pressure, reported from a HFOV &
  Real &
  cm H$_2$O &
  $[0, \infty)$ \\
  $P_{aw}^{PEEP}$ &

  Positive end-expiratory pressure &
  Real &
  cm
  H$_2$O &
  $[0, \infty)$ \\

  \rowcolor{gray!20}
  $VENT$ &
  Indicator for invasive mechanical ventilation &
  Integer &
  &
  $\{0, 1\}$ \\

  $O_{2}Support$ &
  Any non-invasive oxygen support &
  Integer &
  &
  $\{0, 1\}$ \\
\end{longtable}

\subsubsection{Constructed Variables}\label{sec:resp-constructed}

From the EHR values we construct several variables which are then used
for the respiratory score. The constructed variables combine
measurements that are not always recorded at the same time (e.g.,
$SpO_2$ and $FiO_2$) and apply clinical logic about timing and
support. This keeps us from pairing an outdated $FiO_2$ with a more
recent oxygen measurement or vice versa.

\paragraph{PF Ratio}

In general, the PF ratio is defined as
\begin{equation}
  \mathrm{PFR} \coloneqq \frac{PaO_2}{FiO_2}.
\end{equation}

For the determining the PF ratio at time $t$
we consider look-back windows for both $FiO_2$ and $PaO_2$
and the relative look-back between the two. Specifically, the PFR is
only valid when
\begin{equation}
  \tau_{FiO_2}(t) \leq \tau_{PaO_2}(t) \quad \equiv \quad \Delta_{FiO_2}(t) \geq
  \Delta_{PaO_2}(t).
\end{equation}
Conceptually, this restriction is in place to protect from using a
$PaO_2$ value which was taken \emph{before} the most current $FiO_2$
value. That is, PFR is only valid if the $PaO_2$ value is reported at
the same time, or after the $FiO_2$ is reported.

The PFR at time $t$ is:
\begin{equation}
  \PFR \coloneqq
    \begin{cases}
      \frac{\mathcal{C}\left[PaO_2\right]\left(t;m_{resp}\right)}{\mathcal{C}\left[FiO_2\right]\left(t;m_{resp}\right)} & \tau_{FiO_2}(t) \leq \tau_{PaO_2}(t) \\
      \mathrm{NA} & \text{otherwise.}
    \end{cases}
\end{equation}
With this definition, at time $t,$ the PF ratio will either have a
valid numeric value or be $\mathrm{NA}.$

\paragraph{SF Ratio}

In general, the SF ratio is defined as
\begin{equation}
  SFR \coloneqq \frac{SpO_2}{FiO_2}.
\end{equation}

Like the PF ratio, we only consider the SF ratio valid if the $SpO_2$
value is reported at the same time as the $FiO_2$ or after.

\begin{equation}
  \tau_{FiO_2}(t) \leq \tau_{SpO_2}(t) \quad \equiv \quad \Delta_{FiO_2}(t) \geq \Delta_{SpO_2}(t).
\end{equation}

Additionally, the SF ratio is valid if the $SpO_2 \leq 97.$ This requirement is
based on the PODIUM~\citep{bembea2022podium,squires2022podium} organ
dysfunction scoring, and is an additional requirement on the
pSOFA~\citep{matics2017psofa} cut points in
the scoring.

\begin{equation}
  \SFR \coloneqq
  \begin{cases}
    \frac{\mathcal{C}\left[SpO_2\right]\left(t;m_{resp}\right)}{\mathcal{C}\left[FiO_2\right]\left(t;m_{resp}\right)} &
    \tau_{FiO_2}(t) \leq \tau_{SpO_2}(t) \wedge
    \mathcal{C}\left[SpO_2\right]\left(t;m_{resp}\right) \leq 97 \\
    \mathrm{NA} & \text{otherwise.}
  \end{cases}
\end{equation}

\paragraph{Invasive Mechanical Ventilation}

We operationalize the definition for invasive mechanical ventilation
based in part of differences in the EHR records we had access to. In
some sets, a simple VENT indicator variable was provided. In other EHR
sets we had mean airway pressure from either a ventilator or
high-frequency oscillatory ventilation (HFOV). Additionally, we had
positive end-expiratory pressure (PEEP) to consider as well.

We defined invasive mechanical ventilation (IMV) as an indicator
based on four possible conditions.

\begin{equation}
  \IMV \coloneqq \mathbb{I}\left\{\bigvee_{i=1}^{4} A_{i}\left(t; m_{resp}\right)\right\}
\end{equation}
where
\begin{equation}
  \label{eq:imv-conditions}
  \begin{aligned}
    A_{1}\left(t; m_{resp}\right) &\coloneqq \mathcal{C}\left[VENT\right]\left(t;m_{resp}\right) = 1, \\
    A_{2}\left(t; m_{resp}\right) &\coloneqq \mathcal{C}\left[P_{aw}^{VENT}\right]\left(t;m_{resp}\right) > 0, \\
    A_{3}\left(t; m_{resp}\right) &\coloneqq \mathcal{C}\left[P_{aw}^{HFOV}\right]\left(t;m_{resp}\right) > 0, \text{ and} \\
    A_{4}\left(t; m_{resp}\right) &\coloneqq \mathcal{C}\left[P_{aw}^{PEEP}\right]\left(t;m_{resp}\right) > 3.
  \end{aligned}
\end{equation}

The PEEP limit of 3 CM H2O is used to limit false positives while limiting false
negatives.  Similar limits have been used in other criteria.  For example, the
lung injury score\citep{martin2001airway} uses a PEEP of 5 CMH20 for 0 points
and 6-8 CM H20.  Further, \cite{martin2001airway} states a PEEP of 5 CM H20 and
$FiO_2 < 0.5$ as part of the criteria for identifying patients ready to accept a
trial of spontaneous breathing.

\paragraph{Other Respiratory Support}

Extend the definitions in \eqref{eq:imv-conditions} to include
\begin{equation}
  \label{eq:imv-conditions-part2}
  \begin{aligned}
    A_{5}\left(t; m_{resp}\right) &\coloneqq \mathcal{C}\left[FiO_2\right]\left(t;m_{resp}\right) > 0.21, \\
    A_{6}\left(t; m_{resp}\right) &\coloneqq \mathcal{C}\left[O_{2}Support\right]\left(t;m_{resp}\right) = 1. \\
  \end{aligned}
\end{equation}
Then other respiratory support is defined by \eqref{eq:imv-conditions} and \eqref{eq:imv-conditions-part2} as
\begin{equation}
  \label{eq:ors}
  \ORS \coloneqq \mathbb{I}\left\{ \bigvee_{i=1}^{6} A_{i}\left(t; m_{resp} \right) \right\}
\end{equation}

\subsubsection{Respiratory Organ Dysfunction Score}\label{resp-ods}

In the development of Phoenix, we used pSOFA respiratory
scoring\citep{matics2017psofa}. The pSOFA respiratory
scoring does not explicitly require respiratory support to score 1
point, however, for the Phoenix definition the requirement of some respiratory
support was made explicit.  It is extremely unlikely that a patient with SFR or
PFR values sufficiently low enough to score at least one point for respiratory
dysfunction would not have at the very least supplemental oxygen.

\begin{equation}
  \DS{Resp}[\m{resp}] \coloneqq
    \ORS \mathbb{I} \left\{B_1\left(t; m_{resp}\right)\right\}
    +
    \IMV \sum_{i=2}^{3} \mathbb{I}\left\{B_i\left(t; m_{resp}\right)\right\}
\end{equation}
where
\begin{equation}
  \begin{aligned}
    B_1\left(t; m_{resp}\right) &\coloneqq \PFR < 400 \vee \SFR < 292 \\
    B_2\left(t; m_{resp}\right) &\coloneqq \PFR < 200 \vee \SFR < 220 \\
    B_3\left(t; m_{resp}\right) &\coloneqq \PFR < 100 \vee \SFR < 148.
  \end{aligned}
\end{equation}

The domain for the respiratory organ dysfunction score is

\begin{equation}
  \DS{Resp}[\m{resp}] \in \{0, 1, 2, 3\} \quad \forall m_{resp}, t.
\end{equation}

\subsection{Cardiovascular}\label{sec:card}

The cardiovascular score has three components, vasoactive medications,
mean arterial pressure, and lactate. A score of 0, 1, or 2, points is
possible for each of the components resulting in a total score of 0 to 6 points.

\subsubsection{Data from the EHR}\label{sec:card-ehr}

The input data from the EHR needed for the cardiovascular dysfunction
scores are: These include vasoactive medications (as indicators of
pharmacologic cardiovascular support), blood pressure measurements (from
cuff or arterial line), lactate, and age (for age-specific mean arterial
pressure thresholds).

\begin{longtable}{lllll}
  \caption{EHR Inputs for Cardiovascular Organ Dysfunction Scoring}
  \label{tab:card-ehr} \\

  \toprule
  \textbf{Variable} &
  \textbf{Description} &
  \textbf{Type} &
  \textbf{Units} &
  \textbf{Domain} \\
  \midrule
  \endhead

  \bottomrule
  \endlastfoot

  \rowcolor{gray!20}
  $Dobutamine$ & Systemic use of vasoactive medication dobutamine & Integer & -- & $\{0, 1\}$ \\
  $Dopamine$ & Systemic use of vasoactive medication dopamine & Integer & -- & $\{0, 1\}$ \\
  \rowcolor{gray!20}
  $Epinephrine$ & Systemic use of vasoactive medication epinephrine & Integer & -- & $\{0, 1\}$ \\
  $Milrinone$ & Systemic use of vasoactive medication milrinone & Integer & -- & $\{0, 1\}$ \\
  \rowcolor{gray!20}
  $Norepinephrine$ & Systemic use of vasoactive medication norepinephrine & Integer & -- & $\{0, 1\}$ \\
  $Vasopressin$ & Systemic use of vasoactive medication vasopressin & Integer & -- & $\{0, 1\}$ \\

  \rowcolor{gray!20}
  $Lactate$ & Lactate & Real & mmol / L & $[0, \infty)$ \\

  $SBP_c$ & systolic blood pressure reported from a cuff & Real & mmHg & $[0, \infty)$ \\
  \rowcolor{gray!20}
  $DBP_c$ & diastolic blood pressure reported from a cuff & Real & mmHg & $[1, 200]$ \\
  $MAP_c$ & mean arterial pressure reported from a cuff & Real & mmHg & $[0, \infty)$ \\
  \rowcolor{gray!20}
  $SBP_a$ & systolic blood pressure reported from arterial line & Real & mmHg & $[0, \infty)$ \\
  $DBP_a$ & diastolic blood pressure reported from arterial line & Real & mmHg & $[1, 200]$ \\
  \rowcolor{gray!20}
  $MAP_a$ & mean arterial pressure reported from arterial line & Real & mmHg & $[0, \infty)$ \\

  $Age$ & Patient's age & Real & months & $[0, 216)$ \\
\end{longtable}

\subsubsection{Constructed Variable}\label{sec:card-constructed}

We build three components: vasoactive medication use, mean arterial
pressure, and lactate. Each component contributes 0--2 points, and
together they sum to the cardiovascular score. The definitions below
formalize how we combine multiple data sources (e.g., cuff vs arterial
line) and how we handle timing.

\paragraph{Vasoactive Medications}

Each of the six vasoactive medications
\begin{equation}
  \boldsymbol{V} \coloneqq \{Dobutamine, Dopamine, Epinephrine, Milrinone, Norepinephrine, Vasopressin \}
\end{equation}
are to be assessed as 0/1 indicators. Either the patient is receiving systemic
vasoactive medication for the purpose of cardiovascular support or they are not.
Medications such as epinephrine are used in many
different situations for different purposes, it is important that the input data
is an indicator of the use of the medication systemically for cardiovascular
support.

We define the contribution to the cardiovascular score from vasoactive
medications as the count of the number of vasoactive medications received,
regardless of dose, with a cap of 2 points.

\begin{equation}
  \mathcal{D}\left[Vasos\right](t; m_{vaso}) \coloneqq \min\left\{2, \sum_{v \in \boldsymbol{V}}
  \mathbb{I}\left\{\mathcal{C}[v](t;m_{vaso}) = 1\right\} \right\}.
\end{equation}

The domain for the contribution to the cardiovascular score from vasoactive
medications is:
\begin{equation}
  \label{eq:vasos}
  \mathcal{D}\left[Vasos\right](t; m_{vaso}) \in \{0,1,2\} \quad \forall m_{vaso},t
\end{equation}

This was based on the use of the vasoactive inotrope score
(VIS)\citep{haque2015association}.  The VIS used the same six vasoactive
medications and based the score on the dosage.  During the development of the
Phoenix criteria we found some lower resourced sites did not have all six
medications.  As such, patients at these lower resourced sites had a ceiling on
the severity of cardiovascular dysfunction based on VIS.  By ignoring the dosage
and just relying on the fact that the patient was receiving the medication, the
Phoenix criteria was determined to have greater utility across high-, medium-,
and low- resourced sites.

\paragraph{Mean Arterial Pressure (MAP)}

Mean arterial pressure (MAP) is non-trivial to define. We have limits
based on age (in months) and a hierarchy for the preferable source of
information used to define the MAP. In practice, MAP may be recorded
directly or inferred from systolic and diastolic pressures. We
prioritize arterial line measurements when available because they are
more reliable in unstable patients, then fall back to cuff measurements.

The hierarchy for the pressure: use the MAP from an arterial line
($MAP_a$) first. If that is not available, then approximate using systolic and
diastolic values reported from the arterial line.
$SBP_a$ and $DBP_a$. If none of the arterial line values are
available, then use values from a cuff preferentially taking a reported mean
arterial pressure, $MAP_c,$, and lastly, approximate via $SBP_c$ and
$DBP_c.$ That is, let
\begin{equation}
  \begin{aligned}
    \widehat{MAP_{1}}(t; m_{bp}) &\coloneqq \mathcal{C}[MAP_a](t;m_{bp}), \\
    \widehat{MAP_{2}}(t; m_{bp}) &\coloneqq \frac{1}{3}\mathcal{C}[SBP_a](t;m_{bp}) + \frac{2}{3}\mathcal{C}[DBP_a](t;m_{bp}), \\
    \widehat{MAP_{3}}(t; m_{bp}) &\coloneqq \mathcal{C}[MAP_c](t;m_{bp}), \text{ and} \\
    \widehat{MAP_{4}}(t; m_{bp}) &\coloneqq \frac{1}{3}\mathcal{C}[SBP_c](t;m_{bp}) + \frac{2}{3}\mathcal{C}[DBP_c](t;m_{bp}).
  \end{aligned}
\end{equation}
Then
\begin{equation}
  \widehat{\widehat{MAP}}(t; m_{bp}) \coloneqq
  \begin{cases}
    \widehat{MAP_{i^{\ast}}}(t; m_{bp}) &
    \begin{aligned}
      & if \exists i^{\ast} \in \left\{1, 2, 3, 4\right\} \quad \text{where} \\
      & i^{\ast} \coloneqq \min \left\{i \in \left\{1, 2, 3, 4\right\}: \widehat{MAP_{i}}(t; m_{bp}) \ne \mathrm{NA} \right\} \\
    \end{aligned} \\
    \mathrm{NA} & \text{otherwise.}
  \end{cases}
\end{equation}

The limits for the MAP are based on age (in months). Age-based
thresholds reflect pediatric physiology: younger patients have lower
expected MAPs, so the cut-points increase with age. Let

\begin{equation}
  \label{eq:theta1}
  \theta_{1}(age) \coloneqq
    \begin{cases}
      31 & age \in [  0,   1) \\
      39 & age \in [  1,  12) \\
      44 & age \in [ 12,  24) \\
      45 & age \in [ 24,  60) \\
      49 & age \in [ 60, 144) \\
      52 & age \in [144, 216),
    \end{cases}
\end{equation}
and
\begin{equation}
  \label{eq:theta2}
  \theta_{2}(age) \coloneqq
    \begin{cases}
      17 & age \in [  0,   1) \\
      25 & age \in [  1,  12) \\
      31 & age \in [ 12,  24) \\
      32 & age \in [ 24,  60) \\
      36 & age \in [ 60, 144) \\
      38 & age \in [144, 216).
    \end{cases}
\end{equation}

\noindent Then, the part of the cardiovascular score based on MAP is:

\begin{equation}
  \label{eq:map}
  \mathcal{D}\left[MAP\right](t; m_{bp}, age) \coloneqq
    \mathbb{I}\left\{\widehat{\widehat{MAP}}(t; m_{bp}) < \theta_{2}(age)\right\} +
    \mathbb{I}\left\{\widehat{\widehat{MAP}}(t; m_{bp}) < \theta_{1}(age)\right\}.
\end{equation}
The domain for this component of the cardiovascular dysfunction score is
\begin{equation}
  \mathcal{D}\left[MAP\right](t; m_{bp}, age) \in \{0, 1, 2\} \quad \forall m_{bp},t.
\end{equation}

\paragraph{Lactate}\label{sec:lactate}

Lactate is treated as a marker of impaired perfusion and severity of shock.
Higher lactate thresholds contribute more points. The definition for lactate
used in the cardiovascular dysfunction score is:
\begin{equation}
  \label{eq:lactate}
  \mathcal{D}\left[Lactate\right](t; m_{lac}) \coloneqq
    \mathbb{I}\left\{\mathcal{C}[Lactate](t;m_{lac}) \geq  5\right\} +
    \mathbb{I}\left\{\mathcal{C}[Lactate](t;m_{lac}) \geq 11\right\}.
\end{equation}
That is, one point for a lactate value of at least 5 mmol/L; two points for a
lactate value of at least 11 mmol/L. The domain for the contribution of lactate
to the cardiovascular dysfunction score is:
\begin{equation}
  \mathcal{D}\left[Lactate\right](t; m_{lac}) \in \{0, 1, 2\} \quad \forall m_{lac},t.
\end{equation}

\subsubsection{Cardiovascular Dysfunction Score}\label{sec:card-ods}
The cardiovascular dysfunction is based on \eqref{eq:vasos}, \eqref{eq:map}, and
\eqref{eq:lactate}.

\begin{equation}
  \label{eq:card}
  \mathcal{D}[Card](t; m_{vaso}, m_{bp}, m_{lac}, age) \coloneqq
    \mathcal{D}\left[Vasos\right]\left(t; m_{vaso}\right) +
    \mathcal{D}\left[MAP\right]\left(t; m_{bp}, age\right) +
    \mathcal{D}\left[Lactate\right]\left(t; m_{lac}\right).
\end{equation}
The domain for the cardiovascular dysfunction score is
\begin{equation}
  \mathcal{D}[Card](t; m_{vaso}, m_{bp}, m_{lac}, age)
  \in \{0, 1, 2, 3, 4, 5, 6\} \quad \forall \boldsymbol{m},t;
\end{equation}

\subsection{Neurological}\label{sec:neuro}

The neurological score is based on the Glasgow Coma score (GCS) and
whether or not the patient has bilaterally fixed pupils. These capture
level of consciousness and brainstem function, respectively.

\subsubsection{Data from the EHR}\label{sec:neuro-ehr}
\begin{longtable}{lllll}
  \caption{EHR Inputs for Neurological Organ Dysfunction Scoring}
  \label{tab:neuro-ehr} \\

  \toprule
  \textbf{Variable} &
  \textbf{Description} &
  \textbf{Type} &
  \textbf{Units} &
  \textbf{Domain} \\
  \midrule
  \endhead

  \bottomrule
  \endlastfoot
  $GCS_{motor}$ & Glasgow Coma Motor Score & Integer & unitless & $\{1, 2, 3, 4, 5, 6\}$ \\
  $GCS_{eye}$ & Glasgow Coma Eye Score & Integer & unitless & $\{1, 2, 3, 4\}$ \\
  $GCS_{verbal}$ & Glasgow Coma Verbal Score & Integer & unitless & $\{1, 2, 3, 4, 5\}$ \\
  $GCS_{total}$ & Total Glasgow Coma Score & Integer & unitless & $\{3, 4, \ldots, 15\}$ \\
  $Pupils$ & Bilaterally fixed pupils & Integer & & $\{0, 1\}$ \\
\end{longtable}

\subsubsection{Constructed Variable}\label{sec:neuro-constructed}

We derive two variables: a time-aligned GCS total (which may need to be
reconstructed from components) and a binary indicator for fixed pupils.
This helps address common EHR issues where component scores are updated
without the total.

\paragraph{GCS}

In general the GCS total is the sum of the three components, eye,
verbal, and motor. In the development of the Phoenix criteria we found
that there were cases where one or two of the components would be
updated but the other component(s) and the total were not updated. When
this occurred we chose to assume it was because there was no change in
the other component scores. This rule prevents missing totals from
artificially lowering the score when only one component is updated.

\begin{equation}
  \label{eq:gcs}
  GCS(t; m_{gcs}) \coloneqq 
  \begin{cases}
    \mathcal{C}[GCS_{total}](t;m_{gcs}) & \Delta_{GCS_{total}}(t) < \min_{i \in \boldsymbol{G}} \Delta_{GCS_{i}}(t) \\ \\
    \sum_{i \in \boldsymbol{G}} \mathcal{C}[GCS_{i}](t;\infty) & 
    \begin{split}
      & \left[  \Delta_{GCS_{total}}(t) = \mathrm{NA} \vee \min_{i \in \boldsymbol{G}} \Delta_{GCS_{i}}(t) \leq \Delta_{GCS_{total}}(t) \right] \wedge \\
      & \min_{i \in \boldsymbol{G}} \Delta_{GCS_{i}}(t) \leq m_{gcs} 
    \end{split} \\ \\
    \mathrm{NA} & \text{otherwise,}
  \end{cases}
\end{equation}
where
\begin{equation}
  \boldsymbol{G} \coloneqq \{motor,eye,verbal\}.
\end{equation}
The use of $\mathcal{C}[\cdot](t;\infty)$ means ``carry forward the most recent value
with no time limit.'' In practice, this is only applied when each of the
three component scores has been observed at least once. So if the latest
eye, verbal, and motor scores were recorded at different times, we take
each component's most recent value and sum them to form a consistent
total, even if those component timestamps are not identical.

The score contribution from GCS is
\begin{equation}
  \label{eq:gcscontribution}
  \mathcal{D}\left[GCS\right]\left(t;m_{gcs}\right) \coloneqq \mathbb{I}\left\{ GCS\left(t; m_{gcs}\right) \leq 10\right\}
\end{equation}

\paragraph{Pupils}
This is treated as a severe
neurologic finding and can dominate the neurologic score.
The bilaterally fixed pupils are defined as.

\begin{equation}
  \label{eq:pupils}
  \mathcal{D}\left[Pupils\right](t; m_{pupil}) \coloneqq 2 \times \mathbb{I}\left\{\mathcal{C}[Pupils](t;m_{pupil}) = 1\right\}.
\end{equation}

\subsubsection{Neurological Organ Dysfunction Score}\label{sec:neuro-ods}

The neurological component on the score is The score is 0, 1, or 2, where
a fixed-pupil finding yields 2 points.

\begin{equation}
  \mathcal{D}\left[Neuro\right]\left(t; m_{gcs}, m_{pupil}\right) \coloneqq
  \max \left\{ \mathcal{D}\left[GCS\right]\left(t; m_{gcs}\right),
  \mathcal{D}\left[Pupils\right]\left(t; m_{pupil}\right) \right\}.
\end{equation}

\subsection{Coagulation}\label{sec:coag}

The coagulation score uses common laboratory markers of clotting and
fibrinolysis. Each marker contributes at most one point, and the total
is capped at 2. This keeps the score sensitive to abnormalities without
overweighting multiple abnormal labs that can occur together in the same
patient.

For readers without a clinical background: Clotting (coagulation) is the
body's process for stopping bleeding by forming fibrin clots.
Fibrinolysis is the counter-process that breaks down those clots once
they are no longer needed. INR is a standardized measure of how long it
takes blood to clot; higher INR means slower clotting and can reflect
coagulopathy, liver dysfunction, or anticoagulation. D-dimer is a
breakdown product of cross-linked fibrin, so elevated levels suggest
recent or ongoing clot formation and breakdown, though it is not
specific to any single disease.

In severe infection, clotting can be activated even without trauma.
Inflammatory signals and endothelial injury can trigger coagulation
pathways, leading to micro-clot formation in small vessels. At the same
time, fibrinolysis can be overwhelmed or dysregulated. In extreme cases
this becomes disseminated intravascular coagulation (DIC), where
clotting and bleeding risk coexist.

\subsubsection{Data from the EHR}\label{sec:coag-ehr}

\begin{longtable}{lllll}
  \caption{EHR Inputs for Coagulation Organ Dysfunction Scoring}
  \label{tab:coag-ehr} \\

  \toprule
  \textbf{Variable} &
  \textbf{Description} &
  \textbf{Type} &
  \textbf{Units} &
  \textbf{Domain} \\
  \midrule
  \endhead

  \bottomrule
  \endlastfoot
  $Platelets$ & platelets & Real & 1000 / $\mu$L & $[0, \infty)$ \\
  $INR$ & INR & Real & unitless & $[0, \infty)$ \\
  $DDimer$ & D-Dimer & Real & mg/L (FEU) & $[0, \infty)$ \\
  $Fibrinogen$ & Fibrinogen & Real & mg/dL & $[0, \infty)$ \\
\end{longtable}

\subsubsection{Constructed Variable}\label{sec:coag-constructed}

There are no variables to construct.

\subsubsection{Coagulation Organ Dysfunction Score}\label{sec:coag-ods}
Each indicator below tests whether a lab value crosses a clinically
relevant threshold. The score sums how many thresholds are met, then
caps at 2 points.

Define:
\begin{equation}
\begin{aligned}
  \mathcal{D}\left[Platelets\right](t; m_{coag})  &\coloneqq \mathbb{I}\left\{\mathcal{C}[Platelets](t;m_{coag}) < 100\right\} \\
  \mathcal{D}\left[INR\right](t; m_{coag})        &\coloneqq \mathbb{I}\left\{\mathcal{C}[INR](t;m_{coag}) > 1.3\right\} \\
  \mathcal{D}\left[DDimer\right](t; m_{coag})     &\coloneqq \mathbb{I}\left\{\mathcal{C}[DDimer](t;m_{coag}) > 2\right\} \\
  \mathcal{D}\left[Fibrinogen\right](t; m_{coag}) &\coloneqq \mathbb{I}\left\{\mathcal{C}[Fibrinogen](t;m_{coag}) < 100\right\},
\end{aligned}
\end{equation}
and then
\begin{equation}
  \mathcal{D}\left[Coag\right](t; m_{coag}) \coloneqq \min \left\{2, \sum_{x \in \boldsymbol{C}} \mathcal{D}\left[x\right]\left(t; m_{coag}\right) \right\},
\end{equation}
where
\begin{equation}
  \boldsymbol{C} \coloneqq \left\{Platelets, INR, DDimer, Fibrinogen\right\}.
\end{equation}


\subsection{Endocrine}

The endocrine score is part of the Phoenix-8 scoring system.

\subsubsection{Data from the EHR}
There is only one element from the EHR used in the endocrine dysfunction
assessment: glucose, Table~\ref{tab:endo-ehr}.

\begin{longtable}{lllll}
  \caption{EHR Inputs for Endocrine Organ Dysfunction Scoring}
  \label{tab:endo-ehr} \\

  \toprule
  \textbf{Variable} &
  \textbf{Description} &
  \textbf{Type} &
  \textbf{Units} &
  \textbf{Domain} \\
  \midrule
  \endhead

  \bottomrule
  \endlastfoot

  \rowcolor{gray!20}
  $Glucose$ &
  Blood Glucose &
  Real &
  mg/dL &
  $[0, \infty)$ \\
\end{longtable}

\subsubsection{Constructed Variables}
There are no constructed variables to define.

\subsubsection{Endocrine Organ Dysfunction Score}

\begin{equation}
  \mathcal{D}[Endo]\left(t; m_{endo}\right) \coloneqq
  \mathbb{I}\left\{\mathcal{C}[Glucose]\left(t; m_{endo}\right) < 50 \vee \mathcal{C}[Glucose]\left(t; m_{endo}\right) > 150 \right\}
\end{equation}
The domain for the endocrine score is
\begin{equation}
  \mathcal{D}[Endo]\left(t; m_{endo}\right) \in \left\{0, 1\right\}, \quad
  \forall m_{endo}, t.
\end{equation}

\subsection{Immunologic}
The immunologic score is part of the Phoenix-8 scoring system.

\subsubsection{Data from the EHR}
Two values are required from the EHR: Absolute neutrophil count  (ANC) and
absolute lymphocyte count (ALC).  Both are expected in units of 1000 cells per
cubic millimeter, Table~\ref{tab:immu-ehr}.

\begin{longtable}{lllll}
  \caption{EHR Inputs for Immunologic Organ Dysfunction Scoring}
  \label{tab:immu-ehr} \\

  \toprule
  \textbf{Variable} &
  \textbf{Description} &
  \textbf{Type} &
  \textbf{Units} &
  \textbf{Domain} \\
  \midrule
  \endhead

  \bottomrule
  \endlastfoot

  \rowcolor{gray!20}
  $ANC$ &
  Absolute neutrophil count &
  Real &
  1000 cells / mm$^3$ &
  $[0, \infty)$ \\

  $ALC$ &
  Absolute lymphocyte count &
  Real &
  1000 cells / mm$^3$ &
  $[0, \infty)$ \\
\end{longtable}

\subsubsection{Constructed Variables}
There are no constructed variables for the immunologic score.

\subsubsection{Immunologic Organ Dysfunction Score}
Define
\begin{equation}
  \mathcal{D}\left[ALC\right]\left(t; m_{immu}\right) \coloneqq
  \mathbb{I} \left\{ \mathcal{C}\left[ALC\right]\left(t; m_{immu}\right) < 1.00 \right\},
\end{equation}
and
\begin{equation}
  \mathcal{D}\left[ANC\right]\left(t; m_{immu}\right) \coloneqq
  \mathbb{I} \left\{ \mathcal{C}\left[ANC\right]\left(t; m_{immu}\right) < 0.50 \right\}.
\end{equation}
Then
\begin{equation}
  \mathcal{D}[Immu]\left(t; m_{immu}\right) \coloneqq
  \max\left\{
  \mathcal{D}\left[ALC\right]\left(t; m_{immu}\right),
  \mathcal{D}\left[ANC\right]\left(t; m_{immu}\right)
    \right\}.
\end{equation}
The domain for the immunologic score is
\begin{equation}
  \mathcal{D}[Immu]\left(t; m_{immu}\right) \in \left\{0, 1\right\}, \quad
  \forall m_{immu}, t.
\end{equation}

\subsection{Hepatic}
\subsubsection{Data from the EHR}
Two values are required from the EHR: total bilirubin reported in units of mg/dL
and alanine aminotrasferase (ALT) reported in units of IU/L, Table~\ref{tab:hepatic-ehr}.

\begin{longtable}{lllll}
  \caption{EHR Inputs for Hepatic Organ Dysfunction Scoring}
  \label{tab:hepatic-ehr} \\

  \toprule
  \textbf{Variable} &
  \textbf{Description} &
  \textbf{Type} &
  \textbf{Units} &
  \textbf{Domain} \\
  \midrule
  \endhead

  \bottomrule
  \endlastfoot

  \rowcolor{gray!20}
  $Bilirubin$ &
  Total bilirubin &
  Real &
  mg/dL &
  $[0, \infty)$ \\

  $ALT$ &
  Alanine aminotransferase &
  Real &
  IU / L &
  $[0, \infty)$ \\
\end{longtable}

\subsubsection{Constructed Variables}
There are no constructed variables for the hepatic organ dysfunction assessment.

\subsubsection{Hepatic Organ Dysfunction Score}

As with endocrine and immunologic scores, the hepatic score can be a 0 or 1 from
either the bilirubin or ALT values.

Define:
\begin{equation}
  \mathcal{D}\left[Bilirubin\right]\left(t; m_{hepatic}\right) \coloneqq
  \mathbb{I} \left\{\mathcal{C}[Bilirubin]\left(t;m_{hepatic}\right) \geq 4 \right\},
\end{equation}
and
\begin{equation}
  \mathcal{D}\left[ALT\right]\left(t; m_{hepatic}\right) \coloneqq
  \mathbb{I}\left\{
    \mathcal{C}[ALT]\left(t;m_{hepatic}\right) > 102 \right\}.
\end{equation}
Then
\begin{equation}
  \mathcal{D}[Hepatic]\left(t; m_{hepatic}\right) \coloneqq
  \max \left\{
    \mathcal{D}\left[Bilirubin\right]\left(t;m_{hepatic}\right),
    \mathcal{D}\left[ALT\right]\left(t;m_{hepatic}\right)
    \right\}.
\end{equation}
The domain for the hepatic score is
\begin{equation}
  \mathcal{D}[Hepatic]\left(t; m_{hepatic}\right) \in \left\{0, 1\right\}, \quad
  \forall m_{hepatic}, t.
\end{equation}

\subsection{Renal}
The renal score is part of the Phoenix-8 scoring system.

\subsubsection{Data from the EHR}
The patient's age (at admission, in months) and their creatinine levels (in
mg/dL) are required to assess renal dysfunction.  Age is also required for
assessment of cardiovascular dysfunction as part of the Phoenix scoring and thus
the only additional variable to collect to extend to Phoenix-8 from Phoenix for
the renal dysfunction score is the creatinine.
\begin{longtable}{lllll}
  \caption{EHR Inputs for renal Organ Dysfunction Scoring}
  \label{tab:renal-ehr} \\

  \toprule
  \textbf{Variable} &
  \textbf{Description} &
  \textbf{Type} &
  \textbf{Units} &
  \textbf{Domain} \\
  \midrule
  \endhead

  \bottomrule
  \endlastfoot

  \rowcolor{gray!20}
  $Creatinine$ &
  Creatinine &
  Real &
  mg / dL &
  $[0, \infty)$ \\

  $Age$ & Patient's age & Real & months & $[0, 216)$ \\
\end{longtable}

\subsubsection{Constructed Variables}
There are no constructed variables required for the assessment of renal
dysfunction.

\subsubsection{Renal Organ Dysfunction Score}
The renal dysfunction score is based on an age adjusted creatinine level
\eqref{eq:renal-score}.  The cut points for the age ranges are the same as used
in the cardiovascular scoring, \eqref{eq:theta1} and \eqref{eq:theta2}.

\begin{equation}
  \label{eq:renal-score}
  \mathcal{D}[Renal]\left(t; m_{renal}, age\right) \coloneqq
  \begin{cases}
    1 & age \in [  0,   1) \wedge \mathcal{C}[Creatinine]\left(t;m_{renal}\right) \geq 0.8 \\
    1 & age \in [  1,  12) \wedge \mathcal{C}[Creatinine]\left(t;m_{renal}\right) \geq 0.3 \\
    1 & age \in [ 12,  24) \wedge \mathcal{C}[Creatinine]\left(t;m_{renal}\right) \geq 0.4 \\
    1 & age \in [ 24,  60) \wedge \mathcal{C}[Creatinine]\left(t;m_{renal}\right) \geq 0.6 \\
    1 & age \in [ 60, 144) \wedge \mathcal{C}[Creatinine]\left(t;m_{renal}\right) \geq 0.7 \\
    1 & age \in [144, 216) \wedge \mathcal{C}[Creatinine]\left(t;m_{renal}\right) \geq 1.0 \\
    0 & \text{otherwise.}
  \end{cases}
\end{equation}

\subsection{Suspected Infection}\label{sec:suspected-infection}

A patient is defined to have a suspected infection if they have at least
one infectious test \emph{ordered} and at least one dose of a systemic
anti-infective agent. Results of the tests are not needed, just that the
tests have been ordered. Anti-infective agents include anti-bacterials,
anti-fungals, anti-virals, or anti-mycobacterial (anti-buberculosis). This
definition is operational: it identifies when clinicians treated the case as
infectious, not whether infection was ultimately confirmed.

Let $t_{1}$ denote the time when an infectious
test was ordered. Let $t_{2}$ denote the time when the first dose of
a systemic anti-infectious agent was administered.

Suspected infection status is defined as

\begin{equation}
\SI \coloneqq \mathbb{I}\left\{t_{1} \leq t\right\} \times \mathbb{I}\left\{t_{2} \leq t\right\}.
\end{equation}

\section{Phoenix Sepsis Score}\label{define-phoenix-sepsis-score}

Define the look-back windows as the vector
\begin{equation}
  \boldsymbol{m} \coloneqq \left(\m{resp}, \m{vaso}, \m{bp}, \m{lac}, \m{gcs}, \m{pupil}, \m{coag}, \m{endo}, \m{immu}, \m{hepatic}, \m{renal}\right).
\end{equation}

In this definition, the time interval of interest is the same for both SI and
Phoenix Sepsis Score (PSS), but the maxima are taken separately. That is, a
patient can have $\SI=1$ at one time and the maximum organ dysfunction score at
another. The Phoenix Sepsis Score combines organ dysfunction with evidence of
suspected infection (via SI), so it reflects organ dysfunction \emph{in the
presence of} suspected infection.

\begin{equation}
  \label{eq:pss}
  \PSS{d}{} \coloneqq
  \left[ \max_{t \in T} \SI \right] \left[ \max_{t \in T} \sum_{x \in \Omega_{d}} \mathcal{D}[x]\left(t;\boldsymbol{m};age\right) \right], \quad d \in \left\{4, 8\right\},
\end{equation}
where
\begin{equation}
  \Omega_{4} \coloneqq \{Resp, Card, Neuro, Coag\}
\end{equation}
and
\begin{equation}
  \Omega_{8} \coloneqq \Omega_{4} \cup \{Endo, Immu, Hepatic, Renal\}.
\end{equation}

Sepsis is defined as organ-dysfunction in the presence of a suspected
infection.  Specifically, a score of at least $\sigma > 0$ points:
\begin{equation}
  \label{eq:sepsis}
  \Sepsis{} \coloneqq
  \mathbb{I} \left\{\PSS{4}{}\left(T;\boldsymbol{m}, age \right) \geq \sigma \right\}.
\end{equation}
Septic Shock is sepsis with at least $\kappa$ point(s) from cardiovascular dysfunction.
\begin{equation}
  \label{eq:septicshock}
  \SepticShock{} \coloneqq
  \mathbb{I} \left\{
    \max_{t \in T} \SI \times \max_{t \in T} \left( \sum_{x \in \Omega_{4}} \mathcal{D}[x]\left(t;\boldsymbol{m};age\right)
  \mathbb{I}\left\{\mathcal{D}[Card]\left(t;\boldsymbol{m};age\right) \geq
  \kappa
  \right\} \right)
  \geq \sigma \right\},
\end{equation}
where $\sigma \geq \kappa > 0.$
Note that the specific shock definition does not simply consider any
cardiovascular dysfunction in the window $T.$  It must occur with the organ
dysfunction associated with reaching the sepsis score threshold.

For the published Phoenix Sepsis Criteria
\citep{sanchezpinto_2024_development,schlapbach_2024_international} the
following parameter values in Table~\ref{tab:phoenix-parmeters} were used:

\begin{longtable}{lrl}
  \caption{Parameter values defining time windows, look-backs, and thresholds
  for the Phoenix Sepsis Score, Sepsis, and Septic Shock.}
  \label{tab:phoenix-parmeters} \\

  \toprule
  \multicolumn{1}{c}{\textbf{Parameter}} & \multicolumn{1}{c}{\textbf{Value}} & \multicolumn{1}{c}{\textbf{Notes}} \\
  \midrule
  \endfirsthead

  \multicolumn{3}{c}
  {{\bfseries \tablename\ \thetable{} -- continued from previous page}} \\
  \toprule
  \multicolumn{1}{c}{\textbf{Parameter}} & \multicolumn{1}{c}{\textbf{Value}} & \multicolumn{1}{c}{\textbf{Notes}} \\
  \midrule
  \endhead

  \midrule
  \multicolumn{3}{r}{{Continued on next page}} \\
  \midrule
  \endfoot

  \bottomrule
  \endlastfoot

  \addlinespace[0.5em]
  \multicolumn{3}{l}{\textbf{Time Window}} \\
  \rowcolor{gray!20}
  ~~$T$ & $\left[0, 1440\right)$ & The first 24-hours of the encounter \\

  \multicolumn{3}{l}{\textbf{Thresholds}} \\
  \rowcolor{gray!20}
  ~~$\sigma$ & 2 & \\
  ~~$\kappa$ & 1 & \\

  \multicolumn{3}{l}{\textbf{Look-backs}} \\
  \rowcolor{gray!20}
    ~~$m_{resp}$  &  360 minutes & all respiratory EHR variables \\
    ~~$m_{vaso}$  &  720 minutes & vasoactive medications \\
  \rowcolor{gray!20}
    ~~$m_{bp}$   &  360 minutes & blood pressure / MAP \\
    ~~$m_{lac}$   &  360 minutes & lactate \\
  \rowcolor{gray!20}
    ~~$m_{gcs}$   &  720 minutes & GCS components or total \\
    ~~$m_{pupil}$ &  720 minutes & pupils \\
  \rowcolor{gray!20}
    ~~$m_{coag}$  & 1440 minutes & coagulation labs \\
    ~~$m_{endo}$  &  720 minutes & glucose \\
  \rowcolor{gray!20}
    ~~$m_{immu}$  & 1440 minutes & ANC and ALC \\
    ~~$m_{hepatic}$  & 1440 minutes & ALT and Total bilirubin \\
  \rowcolor{gray!20}
    ~~$m_{renal}$  & 1440 minutes & creatinine
\end{longtable}

Importantly, restricting the maximization to $t \in T$ does
\textbf{not} truncate the underlying data. For any $t \in T$, the
carry-forward operator $\mathcal{C}[\cdot](t;m)$ can still use
measurements recorded before $T$ (even at negative times) as long as
they fall within the look-back window $m$. Conceptually, this JAMA
definition is time-aligned: at each time $t$ we sum the organ
dysfunction components evaluated at that same time, then take the
maximum of that sum over $T$. This preserves temporal alignment across
organ systems, so a high score reflects concurrent multi-organ
dysfunction at a single time point.

\section{Exploratory Aggregation Schemes}\label{sec:eas}

There are several exploratory aggregation schemes. Daily and block-based ``worst
value'' aggregation is common in pediatric organ dysfunction scoring (e.g.,
PELOD/qPELOD and pSOFA), which motivates the time-decoupled exploratory
aggregation schemes below. See, for example, daily PELOD scoring using the most
abnormal value each day and pSOFA evaluations that use maximum/mean daily scores
or the highest score in fixed time blocks.
\citep{leteurtre2003pelod,bestati2010pelod,matics2017psofa,badke2025psofa}

\subsection{Organ-Level Maxima}\label{sec:eas1}

The primary definition for PSS in \eqref{eq:pss} looks for the max of the sum of
the organ dysfunction scores.  This exploratory scheme looks at the sum of
the max of each of the organ system dysfunction scores.

\begin{equation}
  \label{eq:pss-olm}
  \mathrm{PSS}_{d}^{\mathrm{OLM}}\left(T;\boldsymbol{m}, age \right) \coloneqq
  \left[\max_{t \in T} \SI \right]
  \left[\sum_{x \in \Omega_{d}} \max_{t \in T} \mathcal{D}[x]\left(t;\boldsymbol{m};age\right)\right],
  \quad d \in \left\{4, 8\right\}.
\end{equation}

$\mathrm{PSS}_{d}^{\mathrm{OLM}}$ decouples time across organ systems. It takes the
maximum score for each organ system independently and then sums those maxima.
This measures the overall burden of dysfunction over the interval, not
necessarily dysfunction occurring at the same time. As a result,
$\mathrm{PSS}_{d}^{\mathrm{OLM}}$ can be larger than the $PSS_{d}$, because peaks
from different time points can accumulate. That is, by construction,
\begin{equation}
  PSS_{d}\left(T; \boldsymbol{m}, age\right) \leq \mathrm{PSS}_{d}^{\mathrm{OLM}}\left(T; \boldsymbol{m}, age\right).
\end{equation}
This follows from the max-sum inequality:
\begin{equation}
  \max_t \sum_x f_x(t) \leq \sum_x \max_t f_x (t).
\end{equation}

Using \eqref{eq:pss-olm} to define sepsis is not structurally different from the
primary definition \eqref{eq:sepsis}.
\begin{equation}
  \label{eq:sepsis-olm}
  \substack{Sepsis \\ OLM}\left(T;\boldsymbol{m},age\right) \coloneqq \mathbb{I}
  \left\{\substack{PSS_{4} \\ OLM}\left(T;\boldsymbol{m}, age \right) \geq \sigma \right\}.
\end{equation}
Septic Shock, for organ-level maxima, is an easier calculation than the primary
definition \eqref{eq:septicshock}.
\begin{equation}
  \label{eq:septicshock-olm}
  \substack{SepticShock \\ OLM}\left(T;\boldsymbol{m},age\right) \coloneqq
  \left[
  \substack{Sepsis \\ OLM}\left(T;\boldsymbol{m},age\right)
  \right]
  \left[
    \mathbb{I} \left\{ \max_{t \in T} \mathcal{D}\left[Card\right]\left(t;\boldsymbol{m},age\right) \geq \kappa\right\}
  \right]
\end{equation}
for $\sigma \geq \kappa > 0.$

\subsection{Organ-Level + Cardiovascular Component Decoupling}

In this experimental aggregation schema we decouple all organ systems
\textbf{and} the three cardiovascular components: vasoactive medications, MAP,
and lactate.

\begin{multline}
  PSS_{d}^{CCD} 
  \left(T; \boldsymbol{m}, age\right) \coloneqq \\
  \left[\max_{t \in T} \SI \right]
  \Bigg[
    \sum_{x \in \Omega_{d}\setminus \left\{Card\right\}} \max_{t \in T} \mathcal{D}[x]\left(t;\boldsymbol{m};age\right) +
    \sum_{x \in \boldsymbol{C}_{card}} \max_{t \in T} \mathcal{D}[x]\left(t;\boldsymbol{m};age\right)
  \Bigg],
  \quad d \in \left\{4, 8\right\},
\end{multline}
where
\begin{equation}
  \boldsymbol{C}_{card} \coloneqq \left\{Vasos, MAP, Lactate\right\}.
\end{equation}
Sepsis and Septic Shock are then defined as:
\begin{equation}
  \label{eq:sepsis-ccd}
  \substack{Sepsis \\ CCD}\left(T;\boldsymbol{m},age\right) \coloneqq \mathbb{I}
  \left\{\substack{PSS_{4} \\ CCD}\left(T;\boldsymbol{m}, age \right) \geq \sigma \right\},
\end{equation}
and
\begin{equation}
  \label{eq:septicshock-ccd}
  \substack{SepticShock \\ CCD}\left(T;\boldsymbol{m},age\right) \coloneqq
  \left[
  \substack{Sepsis \\ CCD}\left(T;\boldsymbol{m},age\right)
  \right]
  \left[
    \mathbb{I} \left\{ \sum_{x \in \boldsymbol{C}_{card}}\max_{t \in T} \mathcal{D}\left[x\right]\left(t;\boldsymbol{m},age\right) \geq \kappa\right\}
  \right]
\end{equation}
for $\sigma \geq \kappa > 0.$



\subsection{Full Component Decoupling}

\emph{All} components within each organ system are taken independently. That is,
each term that contributes to an organ score is maximized or minimized
(whichever is worse) independently over $T$, and then combined using the same
capping rules as the original score.  This requires a specific definition for respiratory scoring:
\begin{equation}
  \begin{split}
    \substack{Resp \\ FCD}(T; m_{resp}) \coloneqq{}
    & \max_{t \in T} \left( \IMV \right) \times \left(\mathbb{I}\{B_3(T; m_{resp})\} + \mathbb{I}\{B_2(T; m_{resp})\}\right) \\
    &+ \max_{t \in T} \left(\ORS\right) \times \mathbb{I}\{B_1(T; m_{resp})\},
  \end{split}
\end{equation} where \begin{equation}
\begin{aligned}
  B_1(T; m_{resp}) &\coloneqq \min_{t \in T} \PFR < 400 \vee \min_{t \in T} \SFR < 292, \\
  B_2(T; m_{resp}) &\coloneqq \min_{t \in T} \PFR < 200 \vee \min_{t \in T} \SFR < 220, \\
  B_3(T; m_{resp}) &\coloneqq \min_{t \in T} \PFR < 100 \vee \min_{t \in T} \SFR < 148.
\end{aligned}
\end{equation}
The overall scores are then defined as:
\begin{multline}
    PSS_{d}^{FCD}
    \left(T, \boldsymbol{m}, age\right) \coloneqq \\
    \shoveleft{ \quad \quad
      \left[\max_{t \in T} \SI \right]
      \Bigg[
        \substack{Resp \\ FCD}\left(T; m_{resp}\right) +
        \min \left\{ 2, \sum_{x \in \boldsymbol{C}_{coag}} \max_{t \in T} \mathcal{D}\left[x\right]\left(t;\boldsymbol{m}\right)\right\}
    } \\
    \quad + \sum_{x \in \omega_{d}} \max_{t \in T} \mathcal{D}\left[x\right]\left(t; \boldsymbol{m}, age\right)
    \Bigg], \quad d \in \left\{4, 8\right\},
\end{multline}
where
\begin{equation}
  \omega_4 \coloneqq \left\{Vasos, MAP, Lactate, Neuro\right\},
\end{equation}
and
\begin{equation}
  \omega_8 \coloneqq \omega_4 \cup \left\{Endo, Immu, Hepatic, Renal\right\}.
\end{equation}
Note: the contributions from neurologic, immunologic, and hepatic scores do
consider more than one input form the EHR.  However, because
\begin{equation}
  \max_{t \in T} \max \left\{
  \mathcal{D} \left[GCS\right]\left(t;\boldsymbol{m}\right),
  \mathcal{D} \left[Pupils\right]\left(t;\boldsymbol{m}\right)
  \right\}
  =
  \max \left\{
  \max_{t \in T} \mathcal{D} \left[GCS\right]\left(t;\boldsymbol{m}\right),
  \max_{t \in T} \mathcal{D} \left[Pupils\right]\left(t;\boldsymbol{m}\right)
  \right\}
\end{equation}
we need not explicitly decouple the inputs in the equations.

Sepsis is defined as:
\begin{equation}
  \label{eq:sepsis-fcd}
  \substack{Sepsis \\ FCD}\left(T;\boldsymbol{m},age\right) \coloneqq \mathbb{I}
  \left\{\substack{PSS_{4} \\ FCD}\left(T;\boldsymbol{m}, age \right) \geq
  \sigma \right\},
\end{equation}
for $\sigma > 0.$

Septic shock is defined as follows:
\begin{equation}
  \label{eq:septicshock-fcd}
  \substack{SepticShock \\ FCD}\left(T;\boldsymbol{m},age\right) \coloneqq
  \left[
  \substack{Sepsis \\ FCD}\left(T;\boldsymbol{m},age\right)
  \right]
  \left[
    \mathbb{I} \left\{ \sum_{x \in \boldsymbol{C}_{card}}\max_{t \in T} \mathcal{D}\left[x\right]\left(t;\boldsymbol{m},age\right) \geq \kappa\right\}
  \right]
\end{equation}
for $\sigma \geq \kappa > 0.$

% \subsection{Alternative XX (Persistence Rule)}\label{alternative-XX-persistence-rule}
% Define a persistence requirement over a duration $\delta>0$

% \subsection{Alternative XX (Component-Level Persistence)}
%
% Alternative XX extends the persistence idea to each component within each
% organ system. That is, each component must meet its threshold
% continuously for at least $\delta$ before contributing to the organ
% score.
%
%
%\subsection{Alternative XX (Component-Level Persistence + Organ Decoupling)}
%
%Alternative XX combines the component-level persistence of Alternative XX
%with full organ-system decoupling (as in Alternative 1).
%
% \subsection{Alternative XX (Infection-Anchored Window)}\label{alternative-XX-infectionanchored-window}
% Define an infection-anchored window


\bibliography{../references}
\end{document}

%%%%%%%%%%%%%%%%%%%%%%%%%%%%%%%%%%%%%%%%%%%%%%%%%%%%%%%%%%%%%%%%%%%%%%%%%%%%%%%%
%                                 End of File                                  %
%%%%%%%%%%%%%%%%%%%%%%%%%%%%%%%%%%%%%%%%%%%%%%%%%%%%%%%%%%%%%%%%%%%%%%%%%%%%%%%%
